%%BLOCK after="Fibrancy passes to the union"
\begin{dcomment}{development}{p46:Ce sera bien sûr une catégorie de Thomason-Kan.}{A small ind-object that is homotopy invariant sees only 1-types}
The functors $H$ of these pages are ind-represented in $\mathrm{Cat}\,\mathrm{hot}$, the category of small categories and functors modulo zigzags of natural transformations, and they are expected to depend only on the homotopy type of their argument. For an ind-object indexed by a small filtered category the two properties together are very restrictive. Write $\Pi_1X$ for the fundamental groupoid of a category $X$ (all arrows inverted); it is the reflection of $X$ into groupoids in $\mathrm{Cat}\,\mathrm{hot}$ (commentary at ms.~37). The proposition below was not found in the sources consulted; no priority is claimed.

\textbf{Theorem (Kan--Thurston with a cardinal bound; not proved here).}
\textit{Let $Y$ be a connected CW complex and $\mu$ a cardinal. There are a group $G$ and a map $t:BG\to Y$ whose homotopy fibres are connected and integrally acyclic, such that $\pi_1(t):G\to\pi_1(Y)$ is surjective with perfect kernel $P$ and every homomorphism from $G$ to a group of cardinality at most $\mu$ is trivial on $P$.}

Without the last clause this is the theorem of Kan and Thurston. The proof of the bounded form (acyclic simple overgroups of prescribed cardinality, and an induction over skeleta in which $BG$ is built as a graph of aspherical spaces with injective edge groups) is not reproduced in this commentary.

\textbf{Lemma.}
\textit{Let $t:X\to Y$ be a map of CW complexes with connected, integrally acyclic homotopy fibres. Then for every space $Z$ the map $[Y,Z]\to[X,Z]$ of free homotopy classes is injective.}

\textit{Proof.}
Let $D$ be the double mapping cylinder of $Y\leftarrow X\to Y$ and $c:D\to Y$ the codiagonal. Since $\pi_1(t)$ is surjective on each component, van Kampen's theorem makes $\pi_1(c)$ an isomorphism; since $t$ induces isomorphisms on homology with all local coefficients pulled back from $Y$, the Mayer--Vietoris sequences give the same for $c$. So $c$ is a weak equivalence. Two maps $Y\to Z$ that become homotopic on $X$ define a map $D\to Z$, which factors up to homotopy through $c$; hence they are homotopic.\hfill$\square$

\textbf{Proposition.}\ours
\textit{Let $I$ be a nonempty small filtered category, $K:I\to\mathrm{Cat}\,\mathrm{hot}$ a diagram of small categories $K_i$, and
\[
 H_K(X)=\varinjlim_i\operatorname{Hom}_{\mathrm{Cat}\,\mathrm{hot}}(X,K_i)=\varinjlim_i\pi_0\underline{\operatorname{Hom}}(X,K_i).
\]
If $H_K$ sends every weak equivalence to a bijection, then it sends every functor inducing an equivalence of fundamental groupoids to a bijection, and the canonical morphism of ind-objects $(K_i)_i\to(\Pi_1K_i)_i$ is an isomorphism in $\operatorname{Ind}(\mathrm{Cat}\,\mathrm{hot})$.}

\textit{Proof.}
The functor $q:\mathrm{Cat}\,\mathrm{hot}\to\mathrm{Hot}$ is the localization by the weak equivalences. By hypothesis $H_K=H'q$ for a presheaf $H'$ on $\mathrm{Hot}$. Left Kan extension along $q$ sends the presheaf represented by $K_i$ to the one represented by $BK_i$, commutes with colimits, and is inverse to restriction on presheaves that factor through $q$; hence $H'(Y)=\varinjlim_i[Y,BK_i]$.

Let $\mu$ be an infinite cardinal bounding the sets of arrows of all the $K_i$; it exists because $I$ is small. Let $Y$ be a CW complex with components $Y_c$, choose $t_c:BG_c\to Y_c$ as in the theorem, and put $T=\coprod_cBG_c$, $t=\coprod_ct_c$, $\tau_1Y=\coprod_cB\pi_1(Y_c)$, with $v:Y\to\tau_1Y$ the canonical map. The composite $vt$ is induced by the surjections $G_c\to\pi_1(Y_c)$. For every $i$, restriction along these surjections is an isomorphism of categories
\[
 \underline{\operatorname{Hom}}\Bigl(\coprod_c\pi_1(Y_c),K_i\Bigr)\longrightarrow\underline{\operatorname{Hom}}\Bigl(\coprod_cG_c,K_i\Bigr)
\]
(groups regarded as categories with one object): a functor from $G_c$ is an object $k$ with a homomorphism $G_c\to\operatorname{Aut}(k)$, which kills the kernel because $\operatorname{Aut}(k)$ has at most $\mu$ elements, and a natural transformation commutes with the action of $G_c$ if and only if it commutes with that of its quotient. Taking components and the colimit over $i$, the composite
\[
 H'(\tau_1Y)\xrightarrow{\ H'(v)\ }H'(Y)\xrightarrow{\ H'(t)\ }H'(T)
\]
is bijective. The second map is injective, as a filtered colimit of the injections $[Y,BK_i]\to[T,BK_i]$ of the lemma. So $H'(v)$ is bijective for every $Y$. A functor $f$ inducing an equivalence of fundamental groupoids has $\tau_1(Bf)$ a homotopy equivalence, and the naturality of $v$ gives that $H'(Bf)$ is bijective.

For the last assertion write $\eta_i:K_i\to\Pi_1K_i$ for the units; they induce $\eta_*:H_K\to H_{\Pi_1K}$. By what was proved, $H_K(\eta_i)$ is bijective; let $\theta_i\in H_K(\Pi_1K_i)$ be the element mapped to the class of $\id_{K_i}$. The $\theta_i$ are compatible and define $\theta:H_{\Pi_1K}\to H_K$. On the class of $\id_{K_i}$ the composite $\theta\eta_*$ is the identity by construction; and $H_{\Pi_1K}(\eta_i)$ is bijective (the functor $H_{\Pi_1K}$ factors through $\Pi_1$) and sends both $\eta_*\theta_i$ and the class of $\id_{\Pi_1K_i}$ to the class of $\eta_i$, so $\eta_*\theta$ is the identity as well. An isomorphism of the presheaves is an isomorphism of the ind-objects.\hfill$\square$

So a functor $H$ on $\mathrm{Cat}\,\mathrm{hot}$ that is invariant under weak equivalences is ind-representable by a \emph{small} filtered system only if it factors through the fundamental groupoid, and then the system may be taken to consist of groupoids. Equivalently: a homotopy invariant functor that distinguishes two spaces with the same fundamental groupoid is not ind-represented in $\mathrm{Cat}\,\mathrm{hot}$ by any small filtered system. The statement concerns ind-objects of the quotient category $\mathrm{Cat}\,\mathrm{hot}$ only; it says nothing about systems of categories $K_n$ kept with their actual functors, as in the replacement construction of the preceding commentaries.
\end{dcomment}
